\chapter{线性注意力的动机}

\section{二次复杂度的根本问题}

标准 Transformer 注意力机制的计算和存储复杂度均为序列长度的二次方。具体而言，注意力分数矩阵 $\mathbf{A} = \text{softmax}(\mathbf{Q}\mathbf{K}^\top / \sqrt{d}) \in \mathbb{R}^{N \times N}$ 的计算和存储是最主要的瓶颈。

我们可以从矩阵乘法的角度来理解这一瓶颈。设 $\mathbf{Q} \in \mathbb{R}^{N \times d}$，$\mathbf{K} \in \mathbb{R}^{N \times d}$，则 $\mathbf{Q}\mathbf{K}^\top$ 的结果是 $N \times N$ 矩阵。即使 $d \ll N$，当 $N$ 很大时，$N^2$ 的规模使得计算变得不可行。

\section{为什么不能简单优化矩阵乘法？}

一个自然的想法是：$\mathbf{Q}\mathbf{K}^\top$ 本质上就是两个矩阵相乘，是否存在更高效的计算方式？答案是：对于一般的稠密矩阵乘法，$\mathcal{O}(N^2 d)$ 是最优的，无法进一步降低复杂度。

但注意，我们的最终目标并不是得到注意力分数矩阵本身，而是计算注意力输出 $\mathbf{A}\mathbf{V}$。如果能够避免显式构造 $N \times N$ 的注意力矩阵，就有可能降低复杂度。

\section{结合律的启发}

考虑矩阵乘法的结合律：

\begin{equation}
    \text{Attention}(\mathbf{Q}, \mathbf{K}, \mathbf{V}) = \underbrace{\text{softmax}\left(\frac{\mathbf{Q}\mathbf{K}^\top}{\sqrt{d}}\right)}_{\mathbf{A} \in \mathbb{R}^{N \times N}}\mathbf{V}
\end{equation}

如果去掉 softmax 的非线性归一化约束，考虑一个简化的注意力形式：

\begin{equation}
    \mathbf{A}\mathbf{V} = (\mathbf{Q}\mathbf{K}^\top)\mathbf{V}
\end{equation}

利用矩阵乘法的结合律，我们可以改变计算顺序：

\begin{equation}
    (\mathbf{Q}\mathbf{K}^\top)\mathbf{V} = \mathbf{Q}(\mathbf{K}^\top\mathbf{V})
\end{equation}

其中：

\begin{itemize}
    \item $\mathbf{K}^\top\mathbf{V} \in \mathbb{R}^{d \times d}$，计算复杂度 $\mathcal{O}(Nd^2)$
    \item $\mathbf{Q}(\mathbf{K}^\top\mathbf{V}) \in \mathbb{R}^{N \times d}$，计算复杂度 $\mathcal{O}(Nd^2)$
    \item 总体复杂度：$\mathcal{O}(Nd^2)$
\end{itemize}

由于 $d \ll N$，$\mathcal{O}(Nd^2)$ 远优于 $\mathcal{O}(N^2 d)$。这就是线性注意力的核心思想——通过改变计算顺序，将复杂度从关于 $N$ 的二次方降低到线性。

\section{Softmax 的挑战}

然而，上述推导忽略了 softmax 函数。softmax 引入了行归一化和指数运算，使得它无法简单地通过结合律来优化。具体来说：

\begin{equation}
    \text{softmax}\left(\frac{\mathbf{Q}\mathbf{K}^\top}{\sqrt{d}}\right)_{ij} = \frac{\exp(\mathbf{q}_i^\top \mathbf{k}_j / \sqrt{d})}{\sum_{k=1}^N \exp(\mathbf{q}_i^\top \mathbf{k}_k / \sqrt{d})}
\end{equation}

分母中的求和依赖于所有键向量 $\mathbf{k}_k$，这导致每一行的归一化因子不同，无法将计算分解为 $\mathbf{Q}(\mathbf{K}^\top\mathbf{V})$ 的形式。

因此，线性注意力的关键问题变为：如何替代 softmax，使得注意力计算能够利用结合律改变计算顺序？

\section{线性注意力的基本思路}

线性注意力通过以下方式解决 softmax 的挑战：

\begin{enumerate}
    \item 使用核函数（kernel function）替代点积相似度 $\mathbf{q}_i^\top \mathbf{k}_j$
    \item 设计可分解的核函数 $\kappa(\mathbf{q}_i, \mathbf{k}_j) = \phi(\mathbf{q}_i)^\top \phi(\mathbf{k}_j)$
    \item 通过特征映射 $\phi$ 将计算改写为：
    \begin{equation}
        \mathbf{O}_i = \frac{\sum_{j=1}^N \kappa(\mathbf{q}_i, \mathbf{k}_j)\mathbf{v}_j}{\sum_{j=1}^N \kappa(\mathbf{q}_i, \mathbf{k}_j)} = \frac{\phi(\mathbf{q}_i)^\top \sum_{j=1}^N \phi(\mathbf{k}_j)\mathbf{v}_j^\top}{\phi(\mathbf{q}_i)^\top \sum_{j=1}^N \phi(\mathbf{k}_j)}
    \end{equation}
    \item 令 $\mathbf{S} = \sum_{j=1}^N \phi(\mathbf{k}_j)\mathbf{v}_j^\top \in \mathbb{R}^{d \times d}$ 和 $\mathbf{z} = \sum_{j=1}^N \phi(\mathbf{k}_j) \in \mathbb{R}^d$，这些聚合量只需一次遍历即可计算
\end{enumerate}

通过这种方式，线性注意力实现了 $\mathcal{O}(Nd^2)$ 的时间复杂度和 $\mathcal{O}(Nd)$ 的空间复杂度。
